\section{Scaling Analysis}
\label{app:scaling}

\subsection{Frozen universe and splits}

The scaling universe contains 30 base graphs and 180 tasks. Candidate-route
counts range from 6 at $m=12$ to 17 at $m=22$; each graph contributes six
distinct feasible-cardinality levels. The global range is
$\phistate=2.3842\times10^{-7}$ to $1.4648\times10^{-3}$, or
$D=2.8342$--6.6227. Phase-0--2 rows are excluded from task generation, model
selection, and coefficient fitting.

The scale-controlled coefficient changes from 172 to 199 because the declared
maximum routing cost changes to $22\times9=198$. The normalized ansatz,
objectives, $p=2$ preparation, $p=3$ continuation, 240-evaluation budget, and
CVaR $\alpha=0.10$ otherwise follow the held-out protocol.

\subsection{Resource guard}

Before optimization, representative fixed-parameter forward passes were timed
at each $m$. A size was permitted only if a numerically valid statevector was
obtained, predicted worker memory was at most 40\% of available RAM, and
predicted optimization time did not exceed 1200 seconds. The guard passed at
$m=12,14,16,18,20$ and failed at $m=22$ because the predicted optimization
time was 2249.3 seconds. The $m=22$ cells were therefore administratively
censored by the preregistered computational budget before optimization. The
observed resource table is retained in the frozen artifact; no result-dependent
code rewrite or GPU extension was allowed.

\subsection{Model selection and response summaries}

Four prespecified response forms were evaluated on development graphs with
leave-one-base-graph-out cross-validation. The one-standard-error rule selected
the simplest linear model defined in the main article for every objective. Development
cross-validated RMSE was 0.33769 for O0, 0.16922 for O2, and 0.19391 for O3.
The graph-level exponent distributions, rather than only their means, are
shown in \cref{fig:scaling}.

At the completed $m=20$ extrapolation split, the preregistered mean $\eta_{O0}$
hypothesis was supported, but the intended O3-versus-O0 objective contrast was
not; its observed sign reversed. The scientific verdict is
\textsc{mixed objective scaling}; $m=22$ supplies no QAOA outcome and no
positive scaling result. Historical Phase-1/2 exponent summaries are read-only
comparisons and never enter the frozen model fit.

\subsection{Optimization adequacy}

Twelve selected development cells compared 240 and 480 objective evaluations.
The prospectively defined high-size optimizer-limited ceiling did not trigger:
no O0/O2/O3 cell at $m=16$ or 18 exceeded the 0.10-decade substantial
$\Delta\gfeas$ threshold. Several $m=12$ and $m=14$ controls were budget-
sensitive, so this check does not prove global optimization. It only prevents
using a prespecified high-size budget artifact as an explanation for the
$m=20$ objective reversal.

The retained execution and censoring census is summarized in
\cref{tab:failure-census}.

\begin{table}[ht]
\centering
\caption{Retained execution and censoring census. A zero denotes a retained planned category with no observed event.}
\label{tab:failure-census}
\small
\begin{tabularx}{\textwidth}{@{}Xrrrr@{}}
\toprule
Stage & Success & \shortstack{Timeout/\\OOM} & \shortstack{Other scientific\\failure} & \shortstack{Resource-\\censored} \\
\midrule
Pilot optimized runs & 504 & 0 & 0 & 0 \\
Held-out $p=2$ preparation & 252 & 0 & 0 & 0 \\
Held-out $p=3$ comparison & 252 & 0 & 0 & 0 \\
Scaling $p=2$ and $p=3$ & 1080 & 0 & 0 & 180 planned cells \\
\bottomrule
\end{tabularx}
\end{table}


The complete retained census---including every zero-count failure category for
the pilot and held-out stages and every split/size/depth/objective cell for the
scaling stage---is included as
\texttt{full\_failure\_census.csv} in Online Resource~2. The compact table above
summarizes its scientific denominators without omitting the detailed public
rows.

\subsection{Censoring interpretation}

Resource censoring is deterministic under the frozen guard and is visible in
all summaries. It is not treated as random missingness, and no imputation,
survival model, or extrapolated $m=22$ QAOA value is reported. A future scaling
study would need a redesigned simulation/resource plan frozen before new
outcomes, rather than an after-the-fact completion of the current series.
The censoring records this study's computational-budget limitation; it is not
a claim that 22 qubits are intrinsically unsimulable, a statevector-memory
limit, or a principled barrier for the quantum algorithm.
